\dykband{EXERCISES IN SWAFL}
{\small Open SWAFL (desktop or online) with the benchmark project. The chapter of the \href{https://github.com/halilunaziru73-creator/Coupled-Water-Fluxes-Pore-Scale-Dynamics-Precision-Irrigation-Mulching-Feedback-Loops/raw/main/Software/v1.1.0/SWAFL_Desktop/docs/manual_v1.1.0/SWAFL_Manual_v1.1.0.pdf}{SWAFL manual} that explains each task is given in brackets; record your values in a table and compare them with what is expected (\emph{Check}).}\par\vspace{3pt}
\begin{exlist}
\item {\color{khblue}[Manual, Chapter 20]} Find your farm on the satellite map, demarcate it and read the vegetation and water indices. Comment on the variation within the field.\par\noindent{\small\itshape Check: the indices vary with soil and crop cover.}
\item {\color{khblue}[Manual, Chapter 20]} Add a water source and let the program propose an irrigation system for the field. Record the proposal and the reasons it gives.\par\noindent{\small\itshape Check: the proposal depends on the slope, the soil and the source.}
\item {\color{khblue}[Manual, Chapter 19]} Import a weather file of your own and switch back to the benchmark record. Compare the yearly rain and the reference evapotranspiration.\par\noindent{\small\itshape Check: the comparison shows how your climate differs.}
\item {\color{khblue}[Manual, Chapter 18]} Enter the field size as a length and a width and as an area, and check that the design layers give the same volumes.\par\noindent{\small\itshape Check: the volumes agree.}
\end{exlist}
